\documentclass[11pt]{article}
\usepackage[a4paper,margin=1in]{geometry}
\usepackage{amsmath,amssymb,amsthm}
\usepackage{booktabs}
\usepackage{graphicx}
\usepackage{xcolor}
\usepackage{hyperref}
\hypersetup{colorlinks=true,linkcolor=blue,urlcolor=blue,citecolor=blue}
\newcommand{\rt}[1]{\texttt{#1}}
\newcommand{\fn}[1]{\texttt{#1}}
\newcommand{\wcc}{\mathrm{wcc}}
\newcommand{\att}{\mathrm{att}}
\newcommand{\ifgraphic}[2]{\IfFileExists{#1}{\includegraphics[width=#2]{#1}}%
  {\fbox{\parbox{#2}{\centering\small figure \texttt{\detokenize{#1}} pending}}}}
\newcommand{\iftable}[1]{\IfFileExists{#1}{\input{#1}}%
  {\fbox{\parbox{0.9\textwidth}{\centering\small table
  \texttt{\detokenize{#1}} pending}}}}

\theoremstyle{plain}
\newtheorem{proposition}{Proposition}

\title{\textbf{Report R18 --- two conservation laws, or frozen domain walls?
Rules 108/201 and 156/198}\\
\large Krylov sector analysis of quantum cellular automata}
\author{Tier 1 analysis}
\date{2026-08-10}

\begin{document}
\maketitle

\begin{abstract}
Hilbert-space fragmentation is often explained by the interaction of two
conservation laws --- charge and dipole moment, or domain walls and
magnetisation (Khemani, Hermele and Nandkishore, arXiv:1912.04300). We test that
account against the four rules R9 places off both axes of its map. The verdict is
that the fragmentation is \textbf{fully constructive from frozen domain walls},
and that the conservation laws, which do exist and are exactly of the promised
type, are \emph{consequences} of the wall grammar rather than an independent
ingredient. In detail: (i) there is no charge at all in the computational basis
--- magnetisation is not conserved, so the charge/dipole ladder never starts;
(ii) there is exactly one conserved domain-wall charge per rule, and the two
families differ in which one --- $W_{108}/W_{201}$ conserve the \emph{staggered}
wall number $S=\sum_i(-1)^i b_i$ while $W_{156}/W_{198}$ conserve the
\emph{total} wall number $D=\sum_i b_i$, with the other broken in each case;
(iii) taking both independent charges together their joint level sets number
$36$ at $N=14$ against $3329$ Krylov sectors, so they are short of the required
cardinality by an exponentially growing factor and cannot be the explanation; (iv) each rule has exactly \emph{one} minimal wall word ---
\rt{00}, \rt{11}, \rt{01}, \rt{10} --- and the \emph{set of its occurrences} is a
complete invariant, verified exhaustively to $N=14$; (v) the wall partition
reproduces not merely the sector count but the entire multiset of sector sizes,
so both of R9's exponents fall out of it, $\rho^2=1.754878$ from an exact
order-3 integer recurrence and $\varphi$ from Fibonacci; and (vi) the conserved
charge is a \emph{function} of the wall set --- for the chiral pair simply
$D=2|W|$ --- so it carries strictly less information than the walls do.
\end{abstract}

\tableofcontents

\section{The four rules and the question}

All four rules are unitary, so by R17 their terminal SCCs are their Krylov
sectors, and all four carry a single Hadamard whose gate fires on one neighbour
pattern:
\begin{center}
\begin{tabular}{llll}
\toprule
rule & tuple & fires when $(x_{i-1},x_{i+1})=$ & family \\
\midrule
$W_{201}$ & \rt{VIII} & $(0,0)$ & the PXP constraint \\
$W_{108}$ & \rt{IIIV} & $(1,1)$ & its spin-flip image \\
$W_{156}$ & \rt{IIVI} & $(1,0)$ & chiral \\
$W_{198}$ & \rt{IVII} & $(0,1)$ & its reflection \\
\bottomrule
\end{tabular}
\end{center}

\noindent R9 reports sector-count bases $\rho^2=1.754878$ for the first pair and
$\varphi$ for the second, with $\rho$ the plastic number. The question is whether
that fragmentation is the shadow of two conserved densities, as in the
dipole-conserving models, or whether it is generated constructively by frozen
walls.

\paragraph{A note on where the tools came from.} The wall grammar and the
charge null-space detector are Tier-2 modules (\fn{qca/walls.py},
\fn{qca/charges.py}); they were run read-only for this report. The wall detector
in \fn{scaling/wall\_charges.py} is an \emph{independent} Tier-1
reimplementation, and a regression test pins its output against the four words
Tier-2 reports, so the central claim rests on two separate implementations
agreeing. The two engines were also checked against each other on the sector
counts and largest-sector sizes of all four rules across $N=6\ldots14$: they
agree term for term.

\section{The conservation laws exist, and they are domain-wall numbers}

Write $b_i = x_i \oplus x_{i+1}$ for the bond (domain-wall) variables, and
\[
D \;=\; \sum_i b_i , \qquad S \;=\; \sum_i (-1)^i b_i .
\]

\begin{proposition}
For $W_{108}$ and $W_{201}$, $S$ is conserved on every sector and $D$ is not.
For $W_{156}$ and $W_{198}$, $D$ is conserved and $S$ is not. This holds at both
\rt{obc0} and \rt{pbc} for every $N$ tested.
\end{proposition}

\noindent Two further facts frame these. First, Tier-2's null-space detector
finds \textbf{no charge at all at range $1$ in the computational basis} for any
of the four rules: single-site magnetisation is not conserved, which is
immediate once one notices that the gate flips a site outright. So the ``charge
plus dipole'' construction has nothing to build on here; the only conserved
densities live in the bond basis. Second, the rank of the certified charge
values on the full $2^N$ space is $3$ including the constant, i.e.\ there are
exactly \textbf{two} independent local conserved quantities per rule --- the
right number for the two-conservation-law story, which is precisely why the
story deserves testing rather than dismissal.

\section{But they are far too coarse to explain the fragmentation}

\begin{center}
\small
\iftable{tab_r18_resolving_obc0.tex}
\end{center}

\noindent Testing the hypothesis fairly means using \emph{both} independent
charges, not just the domain-wall one: the second is the wall count $|W|$, a
range-$2$ charge in the computational basis. For $W_{108}$ their joint level sets
number $10, 15, 21, 28, 36$ over $N=6,8,10,12,14$ --- visibly quadratic --- while
the sectors number $37, 114, 351, 1081, 3329$. The ratio is therefore not merely
large but \emph{exponentially growing}: $3.7, 7.6, 16.7, 38.6, 92.5$.

For $W_{156}/W_{198}$ the two charges coincide, since $D=2|W|$ identically
(\S5), so the table's level count for them is a \emph{lower} bound on what the
symmetry resolves and the second independent charge Tier-2's detector reports is
not identified here. It does not matter for the conclusion: running Tier-2's full
certified charge set on the chiral pair gives $12, 17, 23$ level sets at
$N=8,10,12$ against $48, 124, 323$ sectors, i.e.\ ratios $4.0, 7.3, 14.0$ ---
diverging just the same.

This is the standard diagnostic of fragmentation --- exponentially many Krylov
sectors inside polynomially many symmetry sectors --- and it is exactly what
rules out the two-conservation-law account as an \emph{explanation}. Two charges
of this kind can label $O(N^2)$ things. They cannot label $\rho^{2N}$ things.
Whatever is doing the shattering is not them.

\begin{figure}[htbp]
\centering
\ifgraphic{../../figures/fig_wall_charges_obc0.pdf}{\textwidth}
\caption{\textbf{Walls explain the sectors; the charges cannot.}
(a) Sector count (thick line) against the number of distinct wall-occurrence
sets (hollow markers) --- they coincide exactly at every $N$ --- and against the
number of joint level sets of the conserved charges (dotted), which stays in
single figures. (b) Krylov sectors per symmetry sector: exponentially growing,
so the symmetry cannot be the explanation; the dotted line at $1$ is where a
purely symmetry-resolved model would sit. (c) The dominant root of the exact
integer recurrence satisfied by the number of wall sets, against the growth base
R9 fitted for the sector count --- agreement to six digits, with the exponent
derived rather than fitted.}
\label{fig:walls}
\end{figure}

\section{The wall grammar is enough}

\subsection{One word per rule}

Detection is the Tier-2 criterion: a word $w$ is a wall if, placed on the chain
at either sublattice offset, its sites are invariant under one full cycle for
\emph{every} exterior configuration. Since one cycle has light cone $2$, two
context sites on each side exhaust the exterior dependence and the sixteen
configurations of those four sites are the whole test. Each of the four rules has
exactly one minimal wall word:
\[
W_{108}:\ \rt{00}, \qquad W_{201}:\ \rt{11}, \qquad
W_{156}:\ \rt{01}, \qquad W_{198}:\ \rt{10} .
\]
For $W_{201}$ this is transparent: a site flips only when both its neighbours are
$0$, so two adjacent $1$s can never move, and neither can anything they flank.

\subsection{The occurrence set is a complete invariant}

\begin{proposition}
Two states lie in the same Krylov sector if and only if the rule's wall word
occurs at the same set of positions. Verified exhaustively for all four rules at
\rt{obc0} for $N \le 14$, and at \rt{pbc} for $W_{108}$ and $W_{201}$.
\end{proposition}

\noindent Both halves matter and are reported separately in
\fn{wall\_completeness}: the wall set is \emph{constant} on each sector (it is
conserved) and \emph{injective} across sectors (it separates them). What makes
this a constructive account rather than a restatement is that the label is
computed from a single state, with no reference to the sector it lies in. By
contrast the ``frozen-site signature'' --- the positions and values of the sites
that happen to be constant across a sector --- is also a complete invariant for
all four rules, but it is nearly vacuous, because it is defined \emph{from} the
sector. We record it only to note that it is not the content of this section.

\begin{center}
\small
\iftable{tab_r18_walls_obc0.tex}
\end{center}

\subsection{Both exponents fall out of it}

The number of realisable wall sets satisfies an exact homogeneous linear
recurrence with integer coefficients, found by exact rational elimination rather
than by fitting:
\[
\begin{aligned}
W_{108}, W_{201}:&\quad a(N) = 2a(N-1) - a(N-2) + a(N-3),
  &&\text{dominant root } 1.754878 = \rho^2, \\
W_{156}, W_{198}:&\quad a(N) = a(N-1) + a(N-2),
  &&\text{dominant root } 1.618034 = \varphi .
\end{aligned}
\]
These reproduce R9's fitted sector-count bases to six digits
(Fig.~\ref{fig:walls}c). Moreover the wall partition gives the \emph{whole} size
distribution, not just the count: the multiset of preimage sizes of the map
$x \mapsto W(x)$ equals the multiset of sector sizes at every $N$ tested, so
$D_{\max}$ comes out too --- Fibonacci for $W_{108}/W_{201}$, and R9's
room-packing series $6,9,12,16,20,27,36,48,64,\ldots$ for $W_{156}/W_{198}$.

\subsection{Two details that are not decoration}

\paragraph{The boundary is wall material.} At \rt{obc0} the chain is padded with
a frozen $0$ outside each end, and the wall word must be searched on the padded
string. This is not bookkeeping: $W_{108}$'s word is \rt{00}, the padding
supplies \rt{0}s, and searching the bare chain finds only $616$ of the $1081$
sectors at $N=12$. An earlier pass of this analysis concluded from the bare
search that $W_{156}/W_{198}$ needed the charge \emph{in addition} to the walls;
that was an artefact of the missing padding and is not the case.

\paragraph{One defect, on the ring.} For $W_{156}/W_{198}$ at \rt{pbc} the wall
set is complete except for a single collision at every $N$: the two uniform
states $0^N$ and $1^N$ are both frozen singletons and neither contains \rt{01},
so they share the empty wall set and $n_{\text{sectors}} = n_{\text{wall sets}}
+ 1$ exactly ($48$ vs $47$, $124$ vs $123$, $323$ vs $322$, $844$ vs $843$).
$W_{108}/W_{201}$ have no such defect, because for them $0^N$ and $1^N$ are not
both frozen.

\section{The charge is a shadow of the grammar}

The decisive step is not that the walls suffice but that the charges are
\emph{downstream} of them.

\begin{proposition}
For all four rules, at both boundary conditions and every $N$ tested, the
conserved charge is a function of the wall set. For $W_{156}$ and $W_{198}$ it is
explicitly
\[
D \;=\; 2\,|W| ,
\]
$|W|$ being the number of occurrences of the wall word.
\end{proposition}

\noindent The relation for the chiral pair is transparent once stated: on a ring
the \rt{01} and \rt{10} occurrences are equinumerous, and $D$ counts both, so the
total domain-wall number is exactly twice the number of walls. For
$W_{108}/W_{201}$ the staggered charge $S$ is likewise determined by the wall set
but is not affine in $|W|$ --- it depends on the parities of the wall positions,
which is what ``staggered'' means.

So the implication runs one way only: the wall configuration determines the
charge, the charge does not determine the wall configuration, and the sectors are
the level sets of the former. The conservation law is a coarse-grained readout of
the wall data, not a second ingredient interacting with it.

\section{Answering the question as asked}

\begin{enumerate}
\item \textbf{Is it explained by the interaction of two conservation laws?} No.
There are exactly two independent local conserved quantities, which is the right
number, and one of them is genuinely a domain-wall charge, which is the right
kind. But their joint level sets number $O(N^2)$ --- $36$ of them at $N=14$ ---
against $3329$ Krylov sectors, and the ratio grows from $3.7$ to $92.5$ over
$N=6\ldots14$. They do not have the cardinality to be the explanation.

\item \textbf{Is it fully constructive from the frozen domain walls?} Yes. One
minimal wall word per rule; the set of its occurrences is a complete invariant,
computed locally from a single state; the induced partition reproduces the sector
count, the full size distribution and both of R9's exponents, with the exponents
coming from exact integer recurrences rather than fits. The single exception is
the two uniform states on the ring for the chiral pair.

\item \textbf{And the relation between the two accounts?} The charges are
functions of the wall set --- $D = 2|W|$ for the chiral pair --- so the
constructive account strictly subsumes the symmetry account. This is the
opposite of the dipole-conserving picture, where the charges are given and the
shattering is an emergent extra; here the walls are given and the charges are
the emergent extra.

\item \textbf{An incidental contrast worth keeping.} Which wall-number charge
survives is exactly what separates the two families: the PXP-like pair conserves
the staggered wall number and the chiral pair the total one, and that difference
tracks the difference in their exponents ($\rho^2$ against $\varphi$).
\end{enumerate}

\section{Reproduction}

\begin{verbatim}
python -m qca_fragmentation.scaling.wall_charges --bc obc0 --rebuild
python -m pytest qca_fragmentation/tests/test_wall_charges.py -q
\end{verbatim}

\noindent Everything here comes from \fn{analytics/wall\_charges\_obc0.json}.
The Tier-2 detectors used for the cross-checks in \S2 --- the null-space charge
detector, its gate-local certificate, and the wall grammar --- live in the Tier-2
repository and were not modified; only read. If this analysis is to be adopted
there, \fn{scaling/wall\_charges.py} depends on Tier-1 only for
\fn{core.cycle.succ}, \fn{core.rules} and the union-find in
\fn{graph.wcc.make\_succ}, all of which Tier-2 has equivalents for.

\end{document}
